% ============================================
% MEM-12 (ZH)
% 中文翻译 — SPEC-935-R356
% ============================================
\section*{\textbf{MEM-12} --- 裹尸布撕裂者}
\markboth{MEM-12 --- 裹尸布撕裂者}{}

\noindent \textit{院子里的报丧者}

\noindent \textbf{1945 —— 巴尔巴塞纳收容医院}

\vspace{0.5em}\noindent\rule{\textwidth}{0.5pt}\vspace{0.5em}

\noindent 在收容院里待了二十八年。

\noindent 二十八年的剥落墙壁，二十八年的喊叫成为日常，二十八年的来苏水气味萦绕鼻尖。二十八年没有不隔着铁栅栏的窗户感受过风拂过脸庞。二十八年的孤独，它不治愈——它石化。

\noindent 1945年，我在被沉默者的病区。一个没有人说话的病区，因为所有人都学会了说话没有用。我蜷缩在角落里，日记放在膝上，不停地写。我体内的怪物很安静——在冬眠。我感觉它的重量、它的温热、它缓慢的呼吸。

\noindent 偶尔它在我的胸腔里翻个身，我便感到一阵刺痛。不是疼痛。是\textit{记忆}。它在提醒我它在那里。

\vspace{0.5em}\noindent\rule{\textwidth}{0.5pt}\vspace{0.5em}

\noindent 裹尸布撕裂者在一个冬天出现。

\noindent 1945年6月。寒冷从墙壁的缝隙钻进，不走了。我睡着时双手塞在腋窝里，瑟瑟发抖。其他患者也在抖。我们一起抖，像圈栏里的牲口，等待屠宰。

\noindent 一天夜里，一个声音把我惊醒。

\noindent 不是喊叫。
不是哭泣。
是一声吟唱。

\noindent 纤细。尖锐。像很远的地方一个女人哭泣的声音，时高时低。

\noindent \textit{“裹尸布撕裂者……”}我身旁有人低语。

\noindent 是诺诺老爹。一个从1920年就在收容院的老人——住院二十五年。他用颤抖的手指指着窗户。

\noindent \textit{“她来了。”}

\noindent \textit{“谁？”}

\noindent 他没有回答。只是指着。

\noindent 窗台上，站着一只猫头鹰。

\noindent 它是白色的。不是雪白——是旧骨头的白。羽毛在黑暗中仿佛发光，像它的身体是由凝结的月光做成的。眼睛是黑色的，圆圆的，没有眼皮。它看着我。把头转了180度——一个缓慢、流畅的动作，没有哪只鸟应该做得到。

\noindent \textit{“裹尸布撕裂者，”诺诺老爹重复道，“它一唱，就有人要死。它用喙撕开死者的裹尸布，让灵魂能出来。所以叫这个名字。”}

\noindent 猫头鹰张开了喙。

\noindent 发出的声音不是鸣叫。
是人的尖叫。
一个女人的尖叫——但不是痛苦的。
是\textit{警告}。

\noindent 那尖叫穿透了黑夜。穿透了铁栅栏。穿透了我的胸膛。

\noindent 我体内的怪物\textit{醒了}。

\noindent 我感到它在我身体里抬起头。感到它的眼睛睁开——那双黄眼。它\textit{看见}了裹尸布撕裂者。

\noindent 而裹尸布撕裂者看见了它。

\noindent 猫头鹰歪了歪头。尖叫停了。寂静变得如此绝对，我听见自己的血液在流动。

\noindent 然后裹尸布撕裂者开口了。

\noindent 不是用它的声音。

\noindent 是用怪物的声音。

\noindent 从我体内那个洞窟里出来的声音。

\noindent \textit{“还不是时候。”}

\noindent 猫头鹰展开双翅。巨大。雪白。像一张在黑暗中漂浮的裹尸布。

\noindent \textit{“但快了。”}

\noindent 它飞走了。

\noindent 它的白消失在夜色里。

\vspace{0.5em}\noindent\rule{\textwidth}{0.5pt}\vspace{0.5em}

\noindent 第二天，诺诺老爹被发现死在床上。

\noindent 脸朝着窗户。
嘴张着。
眼睛睁着。

\noindent 眼皮被剥掉了。

\noindent 没有血迹。没有伤痕。眼皮就那么\textit{不见了}。仿佛有人用外科手术般的精准把它们取走了。

\noindent 护工们说是老鼠干的。

\noindent 我知道不是。

\noindent 裹尸布撕裂者取走了诺诺老爹的眼皮。
好让他\textit{再也无法闭上眼睛}。
好让他看见正在到来的东西。
即使在死后。

\vspace{0.5em}\noindent\rule{\textwidth}{0.5pt}\vspace{0.5em}

\noindent 一个星期后，猫头鹰回来了。

\noindent 每天夜里，它落在窗台上看着我。不唱。不叫。就待在那里，一动不动，黑眼睛定定地盯着我。

\noindent 而每天夜里，我体内的怪物开口说话。

\noindent \textit{“它是我的舌头，”它说，“我飞行的眼睛。我在活人世界里的声音。”}

\noindent \textit{“它为什么来？”}

\noindent \textit{“来看看你怎么样了。”}

\noindent \textit{“那我怎么样了？”}

\noindent 沉默。

\noindent \textit{“准备好了。”}

\noindent \textit{“准备好什么？”}

\noindent \textit{“把日记传递下去。”}

\vspace{0.5em}\noindent\rule{\textwidth}{0.5pt}\vspace{0.5em}

\noindent 那一夜之后，我开始用不同的方式书写。

\noindent 笔迹变了。词语来得更快。我不需要思考——手自己动，仿佛有谁在引导着笔。

\noindent 日记不再是我的。

\noindent 怪物正借我之手书写。

\noindent 我只是那只手。
它才是那声音。

\noindent 而它写下的东西——我当时看不懂。后来才懂。当我重读时。当我发现那些词语不是写给我的。

\noindent 是写给后来读到的人。

\noindent 是写给你的。

\vspace{0.5em}\noindent\rule{\textwidth}{0.5pt}\vspace{0.5em}

\noindent 最后一夜，猫头鹰没有来。

\noindent 我在黑暗的牢房里独自醒来，知道它已经走了。

\noindent 但我体内的怪物开口了：

\noindent \textit{“它去铺路了。”}

\noindent \textit{“铺什么路？”}

\noindent \textit{“日记不能留在这里。它必须出去。必须被找到。必须被打开。”}

\noindent \textit{“被谁？”}

\noindent 怪物笑了。笑声在我的胸腔里回荡。

\noindent \textit{“一个像你当初一样需要我的人。”}

\noindent \textit{“那要等到什么时候？”}

\noindent \textit{“等你死的时候。”}

\noindent \textit{“在那之前呢？”}

\noindent \textit{“我们等。”}

\noindent 它安静下来。

\noindent \textit{“我们总是等。”}

\vspace{0.5em}\noindent\rule{\textwidth}{0.5pt}\vspace{0.5em}

\noindent 又是三十四年的等待。

\noindent 三十四年听着裹尸布撕裂者在收容院的屋顶吟唱。

\noindent 三十四年书写着不属于我的词语。

\noindent 三十四年做一个怪物的手。

\noindent 直到1979年，新医生到来。

\noindent 而循环开始闭合。

\vspace{0.5em}
\noindent\textit{\textbf{↪ 路线:} \rota{MEM-11} \rotasep \rota{MEM-13} \rotasep \rota{DOC-07}}
